\section{Nuclear response at all phonon orders}\label{app:D.31}

A coupling that multiplies the density of each species $\kappa$ by a real number $c_\kappa$ probes the crystal through $O(\vect{q}) = \sum_\kappa c_\kappa n_\kappa(\vect{q})$, with $n_\kappa(\vect{q}) = \sum_{\vect{R}}\ee^{\ii\vect{q}\cdot\vect{X}_{\vect{R}\kappa}}$ and $\vect{X}_{\vect{R}\kappa} = \vect{R} + \vect{\tau}_\kappa + \vect{u}_{\vect{R}\kappa}$. From the response of $O(\vect{q})$ in a harmonic crystal in thermal equilibrium, \cref{eq:11.3-exact}, we derive the terms of order $n$ in the phonons, \cref{eq:11.3-Sn}, and the first moment of each order. The average over a zone, the decay of the static correlation, the pointwise band and the bound on the second-order thermal diffuse scattering follow. The notation is that of \cref{sec:5.6}; $N$ is the number of cells, $\Vol = N\Omega$ the volume, and the phonon wavevectors $\vect{Q}$ run over the $N$ points of the first zone.

\subsection{The $n$-phonon term in reciprocal space}

In \cref{eq:5.6-C}, write the absorption term at $-\vect{Q}$ instead of $\vect{Q}$. With $\vect{e}_{\kappa\nu}(-\vect{Q}) = \vect{e}_{\kappa\nu}(\vect{Q})^*$ and frequencies and occupations even in $\vect{Q}$, the emission term and the absorption term take the same form,
\begin{equation}
  \vect{q}\cdot\mathcal{C}_{\kappa\kappa'}(\Delta\vect{R},t)\cdot\vect{q} = \frac1N\sum_{\vect{Q}\nu}\sum_{s = \pm1}\frac{n_{\vect{Q}\nu} + \frac12(1 + s)}{2\omega_{\vect{Q}\nu}\sqrt{M_\kappa M_{\kappa'}}}\;y^s_{\kappa}\,(y^s_{\kappa'})^*\;\ee^{\ii s\vect{Q}\cdot\Delta\vect{R}}\,\ee^{-\ii s\omega_{\vect{Q}\nu}t} ,
  \label{eq:D.31-Cs}
\end{equation}
with $y^{+}_\kappa = \vect{q}\cdot\vect{e}_{\kappa\nu}(\vect{Q})$ and $y^{-}_\kappa = (y^{+}_\kappa)^*$. The term $s = +1$ creates the phonon $\vect{Q}\nu$, the term $s = -1$ absorbs it. The $n$th power of \cref{eq:D.31-Cs}, divided by $n!$, is the term of order $n$ of the exponential in \cref{eq:11.3-exact}. The term of order $n$ is a sum over $n$ phonons $(\vect{Q}_i,\nu_i,s_i)$ with the phase $\ee^{\ii\sum_is_i\vect{Q}_i\cdot\Delta\vect{R}}$. The sum of this phase times $\ee^{-\ii\vect{q}\cdot\Delta\vect{R}}$ over the $N$ lattice vectors gives
\begin{equation*}
  \sum_{\Delta\vect{R}}\ee^{-\ii(\vect{q} - \sum_is_i\vect{Q}_i)\cdot\Delta\vect{R}} = N\sum_{\vect{G}}\delta_{\vect{q} - \sum_is_i\vect{Q}_i,\,\vect{G}} ,
\end{equation*}
the conservation of crystal momentum. The time integral of $\ee^{\ii\omega t}\prod_i\ee^{-\ii s_i\omega_i t}$ is $2\pi\delta(\omega - \sum_is_i\omega_i)$. The factors of the two nuclei collect into one squared modulus, with $(y^{s}_\kappa)^* = \vect{q}\cdot\vect{e}^{s}_{\kappa\nu}$. Counting the factors of $N$, with the Kronecker delta removing one sum over wavevectors, then gives \cref{eq:11.3-Sn}. The factor $1/n!$ comes from the exponential and needs no counting of final states; it also gives the right weight when two phonons coincide, which a sum over unordered final states would have to supply separately. The phase $\ee^{\ii\vect{q}\cdot\vect{\tau}_\kappa}$ carries the unreduced $\vect{q}$, since it came from the static positions, not from the phonons.

\subsection{First moments}

For any function $f(t)$ with $S(\omega) = \Omega^{-1}\int\dd t\,\ee^{\ii\omega t}f(t)$, $\int\dd\omega\,\omega S(\omega) = (2\pi/\Omega)\,\ii f'(0)$. We need the time derivative of \cref{eq:D.31-Cs} at $t = 0$. The thermal factors cancel between $s = +1$ and $s = -1$, as in \cref{eq:5.6-sum}, and
\begin{equation*}
  \partial_t\,\mathcal{C}_{\kappa\kappa'}(\Delta\vect{R},t)\big|_{t=0} = -\frac{\ii}{2\sqrt{M_\kappa M_{\kappa'}}}\,\frac1N\sum_{\vect{Q}}\ee^{\ii\vect{Q}\cdot\Delta\vect{R}}\sum_\nu\vect{e}_{\kappa\nu}(\vect{Q})\,\vect{e}_{\kappa'\nu}(\vect{Q})^\dagger = -\frac{\ii}{2M_\kappa}\,\delta_{\kappa\kappa'}\,\delta_{\Delta\vect{R},\vect{0}}\,\mathbf{1} ,
\end{equation*}
by the completeness of the polarization vectors at each $\vect{Q}$, \cref{eq:5.2-ortho}, and $N^{-1}\sum_{\vect{Q}}\ee^{\ii\vect{Q}\cdot\Delta\vect{R}} = \delta_{\Delta\vect{R},\vect{0}}$. The term of order $n$ of a pair is $f_n(t) = c_\kappa c_{\kappa'}\ee^{-\ii\vect{q}\cdot\Delta\vect{R}_{\kappa\kappa'}}\ee^{-W_\kappa - W_{\kappa'}}[\vect{q}\cdot\mathcal{C}_{\kappa\kappa'}\cdot\vect{q}]^n/n!$, and its derivative at $t = 0$ is proportional to $\vect{q}\cdot\partial_t\mathcal{C}_{\kappa\kappa'}\cdot\vect{q}$, which vanishes for every distinct pair. Every term of a distinct pair thus has zero first moment, at every order and temperature. For the self terms $\vect{q}\cdot\mathcal{C}_{\kappa\kappa}(\vect{0},0)\cdot\vect{q} = 2W_\kappa$ and
\begin{equation*}
  \ii f'_n(0) = c_\kappa^2\,\ee^{-2W_\kappa}\,\frac{n(2W_\kappa)^{n-1}}{n!}\,\frac{q^2}{2M_\kappa} = c_\kappa^2\,E_{R,\kappa}\,\ee^{-2W_\kappa}\,\frac{(2W_\kappa)^{n-1}}{(n-1)!} ,
\end{equation*}
with $E_{R,\kappa} = q^2/2M_\kappa$. Summed over $n$, the Poisson weights add to one and the first moment is $(2\pi/\Omega)\sum_\kappa c_\kappa^2E_{R,\kappa}$, the recoil sum rule \cref{eq:5.6-sum} for the crystal. For one species the share of order $n$ is \cref{eq:11.3-fsum}; with several species the shares of the species are averaged with the weights $c_\kappa^2E_{R,\kappa} \propto c_\kappa^2/M_\kappa$.

\subsection{The average over a zone}

Every term of \cref{eq:11.3-exact} depends on $\vect{q}$ in two ways. The phase $\ee^{-\ii\vect{q}\cdot\Delta\vect{R}_{\kappa\kappa'}}$ changes by order one when $\vect{q}$ moves by $2\pi/|\Delta\vect{R}_{\kappa\kappa'}|$. The vertex is the product of the factors $[\vect{q}\cdot\mathcal{C}_{\kappa\kappa'}\cdot\vect{q}]^n$ and $\ee^{-W_\kappa - W_{\kappa'}}$; it changes under a shift $\delta\vect{q}$ only by relative amounts $\sim 2n\,|\delta\vect{q}|/q$ and $4W_\kappa|\delta\vect{q}|/q$. Average $\vect{q}$ over a cell of reciprocal space with the shape and volume of the first zone, $\vect{q} \to \vect{q} + \vect{p}$ with $\vect{p}$ in the zone, holding the vertex at $\vect{q}$. The phase averages to
\begin{equation}
  j_{\mathrm{BZ}}(\vect{r}) = \frac{\Omega}{(2\pi)^3}\int_{\mathrm{BZ}}\dd^3p\;\ee^{-\ii\vect{p}\cdot\vect{r}} , \qquad \vect{r} = \Delta\vect{R}_{\kappa\kappa'} .
  \label{eq:D.31-jBZ}
\end{equation}
For a lattice vector $\vect{r} = \Delta\vect{R}$ the plane waves $\ee^{-\ii\vect{p}\cdot\Delta\vect{R}}$ are orthogonal on the zone and $j_{\mathrm{BZ}}(\Delta\vect{R}) = \delta_{\Delta\vect{R},\vect{0}}$. In a lattice with one atom per cell every pair is separated by a lattice vector, only the self terms survive, and the zone-averaged response is the incoherent one, \cref{eq:11.3-inc}. Holding the vertex at $\vect{q}$ is the approximation. Letting the vertex vary adds two corrections to the self terms: a correction linear in $\vect{p}$, which averages to zero over the inversion-symmetric zone, and a second-order correction $\sim (nq_{\mathrm{BZ}}/q)^2$. For the distinct terms the linear correction carries $\int_{\mathrm{BZ}}\dd^3p\,\vect{p}\,\ee^{-\ii\vect{p}\cdot\Delta\vect{R}}$, which does not vanish at lattice vectors; the linear correction is $\sim nq_{\mathrm{BZ}}/q$ times the size of the distinct terms. The statement is therefore asymptotic: the zone average of the exact response tends to the incoherent response for $q \gg \max(n, 2W_\kappa)\,q_{\mathrm{BZ}}$.

With a basis, a pair on two sublattices is separated by $\Delta\vect{R} + \vect{\tau}_\kappa - \vect{\tau}_{\kappa'}$, which is never a lattice vector, and $j_{\mathrm{BZ}}$ does not vanish at this separation. Summed over all $\Delta\vect{R}$ it is one, since $\sum_{\Delta\vect{R}}\ee^{-\ii\vect{p}\cdot\Delta\vect{R}} = N\delta_{\vect{p},\vect{0}}$ inside the zone. For a zone replaced by the sphere of equal volume, of radius $q_{\mathrm{BZ}}$, $j_{\mathrm{BZ}}(\vect{r}) = 3j_1(q_{\mathrm{BZ}}r)/q_{\mathrm{BZ}}r$, with $j_1$ the spherical Bessel function; for the bond of silicon, $r = \SI{2.352}{\angstrom}$ and $q_{\mathrm{BZ}}r = 2.68$, $j_{\mathrm{BZ}}$ is 0.44. The four bonds of an atom alone give 1.78, and the farther pairs, with the opposite sign, bring the sum from 1.78 back to one. The cross-sublattice terms therefore survive the zone average, reduced only by the decay of the correlations with distance, and their size has to be computed; \cref{eq:11.3-weight} does so for the $\omega$-integrated response. An average over the directions of $\vect{q}$ at fixed $|\vect{q}|$ acts on each pair as $\langle\ee^{-\ii\vect{q}\cdot\vect{r}}P(\hat{\vect{q}})\rangle$ with $P$ a polynomial of degree $2n$ in $\hat{\vect{q}}$. To evaluate the average for one pair, expand the azimuthal average of $P$ in Legendre polynomials of $\cos\theta$, with $\vect{r}$ along the polar axis, and do the elementary angular integrals. The average then becomes a sum of spherical Bessel functions $j_l(q|\Delta\vect{R}_{\kappa\kappa'}|)$, which fall as $1/qr$. The curves of \cref{fig:11-limit-3a} use this form.

\subsection{The $\omega$-integrated weight}

The convergence of the pair sum \cref{eq:11.3-weight} is set by the acoustic modes at small $\vect{Q}$, which dominate $B_{\kappa\kappa'}$ at large separations. At small $\vect{Q}$, $\vect{e}_{\kappa\nu} \to (M_\kappa/M_{\mathrm{cell}})^{1/2}\hat{\vect{e}}_\nu$, so the factor $\vect{e}\vect{e}^\dagger/\sqrt{M_\kappa M_{\kappa'}}$ becomes $\hat{\vect{e}}_\nu\hat{\vect{e}}_\nu/M_{\mathrm{cell}}$ for every pair, and $N^{-1}\sum_{\vect{Q}} \to \Omega\int\dd^3Q/(2\pi)^3$. At $T = 0$ the weight is $1/2\omega = 1/(2c_\nu Q)$, and since $\int\dd^3Q\,\ee^{\ii\vect{Q}\cdot\vect{r}}/Q = 4\pi/r^2$, the correlation $B_{\kappa\kappa'}$ falls as $\Omega/(M_{\mathrm{cell}}c\,r^2)$ up to a factor of order one that depends on the direction of $\vect{r}$ and the anisotropy of the speeds. The terms of order $n$ fall as $r^{-2n}$ against $4\pi r^2\,\dd r/\Omega$ pairs in a shell, and the sum converges absolutely for $n \ge 2$; at $n = 1$ the sum does not converge, and the phase makes it a distribution in $\vect{q}$, which is evaluated in reciprocal space as \cref{eq:11.1-S1}. At $T > 0$ the weight of the soft modes is $(2n + 1)/2\omega \approx T/(c_\nu Q)^2$, the Fourier transform of $1/Q^2$ is $2\pi^2/r$, and $B_{\kappa\kappa'}$ falls only as $T/r$; at $n = 2$ the shells then contribute $\propto r^2\cdot r^{-2}$ each, and the sum diverges linearly. The divergence is an artifact of summing in real space, and the last subsection removes it.

\subsection{The pointwise band}

The inequality of \cref{eq:11.3-band} is the Cauchy--Schwarz step of \cref{sec:11.3}, with $\mathfrak{P}^{(n)}$ defined as \cref{eq:11.3-Sn} with $\sum_\kappa|a_\kappa|^2$ in place of $|\sum_\kappa a_\kappa|^2$. For the zone average, shift $\vect{q}$ by $\vect{p}$ in the Kronecker delta alone and average $\vect{p}$ over the $N$ wavevectors of the zone. For each set of phonons exactly one $\vect{p}$ satisfies the constraint, so the average replaces the constraint by the factor $1/N$, and the average over $\vect{Q}_1, \ldots, \vect{Q}_{n-1}$ becomes an unconstrained average over all $n$ wavevectors. For each nucleus the summand is then a product of $n$ independent single-phonon factors,
\begin{equation*}
  \frac1N\sum_{\vect{Q}\nu,s}\Bigl[n_{\vect{Q}\nu} + \tfrac12(1 + s)\Bigr]\frac{|\vect{q}\cdot\vect{e}_{\kappa\nu}(\vect{Q})|^2}{2M_\kappa\omega_{\vect{Q}\nu}}\,\delta(\omega_i - s\,\omega_{\vect{Q}\nu}) = 2W_\kappa(\vect{q})\,P^\kappa_1(\hat{\vect{q}},\omega_i) ,
\end{equation*}
convolved in energy, and the convolution of the $n$ factors is $(2W_\kappa)^nP^\kappa_n$. With $c_\kappa^2\ee^{-2W_\kappa}(2\pi/\Omega)/n!$ in front, the convolution gives the term of order $n$ of \cref{eq:11.3-inc}, the equality of \cref{eq:11.3-band}. The self density keeps the momentum constraint of the exact response but drops the interference between the nuclei of a cell; the zone average removes the constraint.

\subsection{Second-order thermal diffuse scattering}

Summing \cref{eq:11.3-Sn} at $n = 2$ over $s_1$ and $s_2$ and integrating over $\omega$, we relabel $\vect{Q}_i \to -\vect{Q}_i$ in every absorbed phonon. The factor $\vect{q}\cdot\vect{e}^{-}(-\vect{Q}_i) = \vect{q}\cdot\vect{e}(\vect{Q}_i)^*$ of an absorbed phonon is then the factor of an emitted phonon at $\vect{Q}_i$, the constraint becomes $\vect{Q}_1 + \vect{Q}_2 = \vect{q} - \vect{G}$ for all four sign choices, and the thermal factors add to $2n + 1$:
\begin{equation}
  \mathfrak{W}^{(2)}(\vect{q}) = \frac{\Omega}{2}\int_{\mathrm{BZ}}\frac{\dd^3Q_1}{(2\pi)^3}\sum_{\nu_1\nu_2}\frac{(2n_1 + 1)(2n_2 + 1)}{4\omega_1\omega_2}\,\Bigl|\sum_\kappa c_\kappa\,\ee^{-W_\kappa}\,\ee^{\ii\vect{q}\cdot\vect{\tau}_\kappa}\,\frac{(\vect{q}\cdot\vect{e}^*_{\kappa\nu_1})(\vect{q}\cdot\vect{e}^*_{\kappa\nu_2})}{M_\kappa}\Bigr|^2 ,
  \label{eq:D.31-W2}
\end{equation}
with $\vect{Q}_2 = \vect{q} - \vect{G} - \vect{Q}_1$ and the labels $1, 2$ for $(\vect{Q}_1,\nu_1)$ and $(\vect{Q}_2,\nu_2)$. \Cref{eq:D.31-W2} is the momentum-space form of \cref{eq:11.3-weight} at $n = 2$, the two forms being related by Poisson summation, and the integral converges at any temperature: each factor $(2n + 1)/2\omega$ grows at most as $T/(cQ)^2$, integrable in three dimensions. The integrand is singular only where both factors are singular together, $\vect{Q}_1 \to 0$ and $\vect{Q}_2 \to 0$, which needs $\vect{k} = \vect{q} - \vect{G} \to 0$.

Split the integral at a radius $Q_s$ below which the acoustic branches are linear. Where $|\vect{Q}_1| > Q_s$ or $|\vect{Q}_2| > Q_s$, at most one factor is singular, and the integral is bounded uniformly in $\vect{k}$ and continuous at $\vect{k} = 0$; these are the regular terms. Where both wavevectors are below $Q_s$, only acoustic branches matter, $\vect{e}_{\kappa\nu} = (M_\kappa/M_{\mathrm{cell}})^{1/2}\hat{\vect{e}}_\nu$, and the amplitude becomes $(\vect{q}\cdot\hat{\vect{e}}_1)(\vect{q}\cdot\hat{\vect{e}}_2)/M_{\mathrm{cell}}$ times $\sum_\kappa c_\kappa\ee^{-W_\kappa}\ee^{\ii\vect{G}\cdot\vect{\tau}_\kappa}$, up to a factor $1 + O(ka)$ from $\ee^{\ii\vect{k}\cdot\vect{\tau}_\kappa}$. With $\coth x \le 1/x + x/3$, the thermal factors are $(2n + 1)/2\omega \le T/\omega^2 + 1/12T$. The singular part of the product is the term of second order in $T/\omega$, $T^2/\omega_1^2\omega_2^2$. With $\omega_\nu \ge c_{\mathrm{min}}Q$ and the completeness of the three acoustic polarizations, $\sum_\nu(\vect{q}\cdot\hat{\vect{e}}_\nu)^2/\omega_\nu^2 \le q^2/(c_{\mathrm{min}}Q)^2$ for each phonon. The remaining integral, extended over all space, follows from the convolution theorem with the transforms $1/4\pi r$ of $1/Q^2$ and $2\pi^2/k$ of $1/r^2$: $\int\dd^3Q\,Q^{-2}|\vect{k} - \vect{Q}|^{-2} = \pi^3/k$. Collecting $\Omega/2$, $1/(2\pi)^3$, $T^2q^4/(M_{\mathrm{cell}}^2c_{\mathrm{min}}^4)$ and $\pi^3/k$ gives \cref{eq:11.3-tds}. The cross terms of first order in $T/\omega$ carry $T/\omega_1^2$ times $1/\omega_2$, whose integral grows only as $\ln(Q_s/k)$, and the zero-point term is finite at $\vect{k} = 0$. In the classical limit with no cut the singular part is exactly $1/k$ times $|\sum_\kappa c_\kappa\ee^{-W_\kappa}\ee^{\ii\vect{G}\cdot\vect{\tau}_\kappa}|^2\Omega T^2J/(16\pi^3M_{\mathrm{cell}}^2)$, with $J = k\int\dd^3Q_1\,h(\hat{\vect{Q}}_1)h(\hat{\vect{Q}}_2)/(Q_1^2Q_2^2)$ and $h = \sum_\nu(\vect{q}\cdot\hat{\vect{e}}_\nu)^2/c_\nu^2$, a number that depends only on the directions of $\vect{k}$ and $\vect{q}$; the bound replaces $h$ by $q^2/c_{\mathrm{min}}^2$, for which $J = \pi^3q^4/c_{\mathrm{min}}^4$.
